% !TEX TS-program = pdflatex
% !TEX encoding = UTF-8 Unicode

% This is a simple template for a LaTeX document using the "article" class.
% See "book", "report", "letter" for other types of document.

\documentclass[11pt, titlepage]{article} % use larger type; default would be 10pt

\usepackage[utf8]{inputenc} % set input encoding (not needed with XeLaTeX)

%%% Examples of Article customizations
% These packages are optional, depending whether you want the features they provide.
% See the LaTeX Companion or other references for full information.

%%% PAGE DIMENSIONS
\usepackage{geometry} % to change the page dimensions
\geometry{a4paper} % or letterpaper (US) or a5paper or....
\geometry{margin=2cm, headsep=5mm, includefoot, includehead}

\usepackage{graphicx} % support the \includegraphics command and options
\graphicspath{{figures/}} % Location of the graphics files
% \usepackage[parfill]{parskip} % Activate to begin paragraphs with an empty line rather than an indent

%%% PACKAGES
\usepackage{booktabs} % for much better looking tables
\usepackage{array} % for better arrays (eg matrices) in maths
\usepackage{paralist} % very flexible & customisable lists (eg. enumerate/itemize, etc.)
\usepackage{verbatim} % adds environment for commenting out blocks of text & for better verbatim
\usepackage{subfig} % make it possible to include more than one captioned figure/table in a single float
% These packages are all incorporated in the memoir class to one degree or another...

%%% HEADERS & FOOTERS
\usepackage{fancyhdr} % This should be set AFTER setting up the page geometry
\pagestyle{fancy} % options: empty , plain , fancy
\fancyhead{}
%%% SECTION TITLE APPEARANCE
\usepackage{sectsty}
\allsectionsfont{\sffamily\mdseries\upshape} % (See the fntguide.pdf for font help)
% (This matches ConTeXt defaults)

%%% ToC (table of contents) APPEARANCE
\usepackage[nottoc,notlof,notlot]{tocbibind} % Put the bibliography in the ToC
\usepackage[titles,subfigure]{tocloft} % Alter the style of the Table of Contents
\renewcommand{\cftsecfont}{\rmfamily\mdseries\upshape}
\renewcommand{\cftsecpagefont}{\rmfamily\mdseries\upshape} % No bold!

\usepackage{lastpage}
\usepackage{rotating}
\usepackage{textcomp} %get the correct micro sec display
\usepackage{float}
%---------- Enable IEEEtran.bst configurations ------
\usepackage{IEEEtrantools}
\usepackage{verbatim}
%----------------------------------------------------

%%% END Article customizations

%%% The "real" document content comes below...

%%% Header %%%%%%%%%%%%%%%%%%%%%%%%%%%%%
\setlength{\headheight}{53pt}
\lhead{EH2760 Management of Projects \\		 
       ESS-NW/ESS-CAR \\
       Leon Fernandez, leonfe@kth.se}
\rhead{Final Report \\
       Version 1 \\
       \thepage(\pageref{LastPage})}
\renewcommand{\headrulewidth}{1pt}
%%%%%%%%%%%%%%%%%%%%%%%%%%%%%%%%%%%%%%%%

%%% Footer %%%%%%%%%%%%%%%%%%%%%%%%%%%%%
%\cfoot{blablablabla}
%\renewcommand{\footrulewidth}{0.4pt}
%%%%%%%%%%%%%%%%%%%%%%%%%%%%%%%%%%%%%%%%
%\date{} % Activate to display a given date or no date (if empty),
         % otherwise the current date is printed 

\begin{document}
%\maketitle
\bstctlcite{BSTcontrol} % IEEEtran.bst controls enabled


%----------------------------------------------------------------------------------------
%	TITLE PAGE
%----------------------------------------------------------------------------------------

\begin{titlepage} % Suppresses displaying the page number on the title page and the subsequent page counts as page 1
	\newcommand{\HRule}{\rule{\linewidth}{0.5mm}} % Defines a new command for horizontal lines, change thickness here
	
	\center % Centre everything on the page
	
	%------------------------------------------------
	%	Headings
	%------------------------------------------------
	
	\begin{figure}
   		\centering
    	\includegraphics[scale=1]{kthLogo.png}
	\end{figure}
	
	\textsc{\LARGE KTH Management of Projects}\\[1cm] % Main heading such as the name of your university/college
	
	\textsc{\Large EH2760, HT 2018}\\[0.5cm] % Major heading such as course name
	
	\textsc{\Large FINAL REPORT}\\[0.5cm] % Major heading as well
	
	%------------------------------------------------
	%	Title
	%------------------------------------------------
	
	\HRule\\[0.4cm]
	
	{\huge\bfseries ESS-NW/ESS-CAR}\\[0.4cm] % Title of your document
	
	\HRule\\[1.5cm]	
	
	%------------------------------------------------
	%	Author(s)
	%------------------------------------------------
	
	\begin{minipage}{0.4\textwidth}
		\begin{flushleft}
			\large
                        \textsc{Jonas Ekman}
			\\
			\textsc{Yini Gao}
                        \\
                        \textsc{Jacob Kimblad}
		\end{flushleft}
	\end{minipage}
	~
	\begin{minipage}{0.4\textwidth}
		\begin{flushright}
			\large
                        \textsc{Leon Fernandez}
			\\
			\textsc{Fredrik Hyyrynen}
                        \\
                        \textsc{Yifan Ruan}
		\end{flushright}
	\end{minipage}
	
	% If you don't want a supervisor, uncomment the two lines below
        % and comment the code above
	%{\large\textit{Author}}\\
	%John \textsc{Smith} % Your name
	\vskip 8cm
	\begin{figure}[H]
   		\centering
    	\includegraphics[scale=0.4]{funLogo.png}
	\end{figure}
	%------------------------------------------------
	%	Date
	%------------------------------------------------
	
	\vfill\vfill\vfill % Position the date 3/4 down the remaining page
	
	{\large\today} % Date, change the \today to a set date if you want
                       % to be precise
	
	%------------------------------------------------
	%	Logo
	%------------------------------------------------
	
	%\vfill\vfill
	%\includegraphics[width=0.2\textwidth]{placeholder.jpg}\\[1cm]
        % Include a department/university logo
        % - this will require the graphicx package
	 
	%-------------------------------------------------------------------
	
	\vfill % Push the date up 1/4 of the remaining page
	
\end{titlepage}

%-------------------------------------------------------------------------

\clearpage

\section{General Summary}
This paragraph should provide the reader with a brief description of how the project was performed
\begin{itemize}
    \item Which results were reached?
    \item Differences between planned outcomes (time, budget, results) and the actual outcome
    \item Discuss the problems that occured and their impact on the project
\end{itemize}


\section{Follow-up of objectives}
\subsection{Goals}
\begin{enumerate}
\item The computer vision system of the car does not have the required performance to
properly follow a green circle held in front of the car.

\item The hardware has been extended to include one unique processor for each sensor or actuator
node and the car has all the functionality it needs to keep a distance of 30 cm to an
object in front.

\item The car has been made failsafe to the extent that it can detect if any of the major
nodes are missing. Most importantly, the node that controls the motor can by itself
stop the car within the required time, even if it gets isolated from the other nodes.
No battery level reading has been implemented.

\item The software running on the controller node checks the availability of all subsystem
and does not start if any of them are missing.

\item The physical layer har been set in a topology that allows the SDN-system
to route packets to and from all nodes.

\item The SDN system consists of a SDN-controller that runs on a Raspberry Pi and Zodiac FX OpenFlow
switches.

\item There is only one route between any two given nodes that are connected as hosts to the
SDN system. As long as the SDN system is activated, packets are always take this route.

\item The SDN controller does not reassign network resources during runtime.
\end{enumerate}

\subsection{Action Plan}
Goals 1, 3 and 8 cannot be considered completed. No action will be taken for goal 3 since
the current situation was deemed acceptable by the stakeholders. The cause of failure for
goal 1 and 8 have been identified and the shareholders have been notified. The action taken
for these failures is to hand over a well-documented repository to the stakeholders so that
they can have another group of students solve these issues in a future course round.

\section{Lessons learned and suggestions for improvement}
The most important part of the final report! The purpose of the section is to reflect on the experience you have had in the project and documenting it
\begin{itemize}
    \item A good way to get a balance picture of the whole project is to think of 5 good and 5 bad experiences during the project.
    \item What went well in the project and why?
    \item What could have been done better?
\end{itemize}

This section should also bring up other factors that have influenced the project such as: seminars, tools, suppliers, tutors, etc.


\section{Summary of time and resource plans}
Comments on the time and resource plan and its milestones and tollgates
\begin{itemize}
    \item Which tasks were dealt with on time?
    \item Which tasks took too long and what was forgotten in the planning?
    \item The final versions of the time and resource plans are attached as an appedix
\end{itemize}


\section{Final comment/other}
This section can include the project manager's sub-project managers' own views and comments on the project, for example: valid documentation, work practices, project management and more.


\section{Attachments}
Final resource plan
Final milestone char (or Gantt Chart if used)


\clearpage

\bibliographystyle{IEEEtran}

\bibliography{reference}
\clearpage
%-------------------------------------------------------------------------
\clearpage
%TODO appendix
\appendix
\clearpage

\end{document}
